Using a strategy similar to the one in section \ref{repeat-token}, we use the $\PE$ to flag when a certain position is reached.
Then, if the transformer is normalized, modifying the last vector we decide the outputted distribution.
Particularly, we can force it to concentrate the probability in $\qaccept$ or $\qreject$.
We change the last vector with a constant function.
The flag inserted by the $\PE$ is to decide if the function has to be applied or not.

\medskip

Notice that this primitive, together with the one described in section \ref{repeat-token}, proves that for every transformer $\TF$ in $\CoTtndn$ there is a transformer $\TF'$ in the same class that always consumes all of its budget.

We extend the alphabet with two dummy versions of the decision tokens and construct $\TF'$ over $\TF$.
We make $\TF'$ work as $\TF$, but we make it print the dummy versions of the decision tokens when $\TF$ would have printed the real ones\footnote{This is achieved giving the probability of each decision token to its dummy version.}.
Then, we force $\TF'$ to repeat the dummy decision tokens until it reaches the last position of the tape (uses all the budget).
At that point, we force $\TF'$ to print the real decision token corresponding to the one it was repeating.